% Dark-profile regression: surfaces, code, citations and restoration after covers.
\documentclass[theme=tokyo-night,language=english,cjk=false]{campusbeamer}
\title{Tokyo Night}
\subtitle{Dark profile regression}
\author{Campus Beamer}
\date{}
\begin{document}
\maketitle
\section{Dark surfaces}
\subsection{Palette and contrast}
\sectiontocpage[overview]
\sectioncoverpage[primary]{The default primary-color section page.}
\sectioncoverpage[white]{An explicitly requested white split remains available.}
\begin{frame}{Semantic callouts}
  \begin{columns}[T,onlytextwidth]
    \begin{column}{.48\textwidth}
      \begin{tipbox}A practical suggestion.\end{tipbox}
      \begin{notebox}Supporting context.\end{notebox}
      \begin{alertbox}An important condition.\end{alertbox}
    \end{column}
    \begin{column}{.48\textwidth}
      \begin{examplebox}A concrete example.\end{examplebox}
      \begin{definitionbox}A precise definition.\end{definitionbox}
    \end{column}
  \end{columns}
\end{frame}
\begin{frame}[campusfocus,campusbackground=theme/tokyo-night/assets/backdrop.pdf,
  campusbackgroundopacity=.2]{Focus on the message}
  Light text remains readable over the dark decorative background.
\end{frame}
\begin{frame}[fragile]{Code and its caption}
\begin{campuscode}[language=Python,caption={A minimal greeting.}]{hello.py}
def greet(name):
    return "Hello " + name
\end{campuscode}
\end{frame}
\begin{frame}{Citations and project artwork}
  \campusdoccite{Project}{https://github.com/AnikiFan/campus-beamer}
  See the project documentation for the full example.
  \fitfigure[width=.55,height=.35]{\schoolwordmarkonlight}{Project lettering on a dark surface}
\end{frame}
\backmatter
\end{document}
